\section{Methods of result visualization}

For the visualization, we chose three methods:
\begin{itemize}
\item Heatmaps: They are useful for visualizing the similarity between users or animes based on their ratings. They will allow us to identify clusters of similar users or animes and explore how their preferences are related to one another.

% https://rpubs.com/jeknov/movieRec
\begin{figure}
\centering
    \caption{An exemplary heatmap}
    \label{fig:heatmap}
\end{figure}

\item Histograms: They will be used to visualize the distribution of ratings and identifying any biases or patterns in the data. They can help us identify whether the ratings are normally distributed or skewed.

\begin{figure}
\centering
    \caption{An exemplary histogram}
    \label{fig:precission-recall}
\end{figure}

\item Precision-recall curves: They show the tradeoff between precision and recall at different thresholds, allowing us to identify the optimal threshold for making recommendations. By visualizing the performance of the algorithm in this way, we can identify areas for improvement and generate new hypotheses for increasing the accuracy of the recommendations.
\end{itemize}

\begin{figure}
\centering
    \caption{An exemplary precision-recall graph}
    \label{fig:histogram}
\end{figure}

